\specialchapter{Introduction}

Spectral theory aims to study the properties of certain normed algebras, the fundamental examples being the algebra of complex-valued continuous functions on a topological space, the algebra of endomorphisms of a Banach space, and the convolution algebra of functions integrable with respect to a Haar measure on a locally compact group. The scope of this study is not necessarily apparent from this simple indication, but to grasp its importance one need only observe that two contemporary mathematicians, unless they belong to the same school, will often have as common knowledge, besides linear algebra and the elementary part of differential and integral calculus, only certain aspects of Fourier analysis.

In the first two chapters of this Book, we develop the most fundamental ideas of spectral theories, and consider in greater detail the third example mentioned above, namely, the Fourier transform, in the setting of a locally compact abelian group.

The following chapters will first continue the general theory of compact linear maps, then that of operators on Hilbert spaces, this case including linear maps defined only on a subspace (generally dense) of such a space. As is well known, this theory is the basis of the mathematical formalism of Quantum Mechanics. Finally, we will present the elementary theory of unitary representations of topological groups, particularly compact ones, which we have already had occasion to use in Chapter IX of the book Lie Groups and Lie Algebras.

The titles of the first two chapters are the same as those of the first edition of this Book, published in 1967. They have been revised, corrected, and expanded.

For ease of reading, we indicate below the most important changes found in this edition, and in particular the additions.

For Chapter I:

- In no. 3 of §3, we introduce the notion of proper partial map, which allows the Gelfand correspondence between locally compact spaces and commutative C*-algebras to be stated in a functorial manner.

- In §4, we added nos. 13 and 14, which deal with holomorphic functional calculus in real or complex normable algebras, and in algebras without a unit element.

- In §6, the presentation of the continuous functional calculus has been reworked (no. 6). Nos. 9 to 11 are new; they consider the notion of a positive element in a general C*-algebra, approximate units of such algebras, and the quotient of a C*-algebra by a closed two-sided ideal.

- §7, concerning the most elementary spectral properties of endomorphisms of Banach spaces, is essentially new; numbers 7 and 8 partially take up certain numbers of the old paragraph 6. It is this new paragraph that will serve as the basis for the more in-depth study of endomorphisms of Hilbert spaces in Chapter IV.

For Chapter II:

- §1 on the Fourier transform has been revised in detail, for example no. 6, which deals with the functoriality of Pontryagin duality. In the new no. 9, we make more explicit the statements most important for the groups \(\mathbf{R}^n\) and \((\mathbf{R}/\mathbf{Z})^n\). Fourier theory on \(\mathbf{R}^n\), whose most natural framework lies in the theory of tempered distributions, is stated here through its essential properties, without giving the proofs, which will appear in Chapter IV.

- We have finally added a Fourier theory summary sheet that brings together the definitions and properties of the main function spaces involved in this theory, along with the essential formulas.

Many of the exercises are new. In particular, we note the considerable increase in the number of exercises in §1 of Chapter II, reflecting the mathematical importance of Fourier theory.

